\section[Répartition spatiale des composantes cinématiquement froide et chaude]{Répartition spatiale des composantes cinématiquement froide et chaude}
\label{sec:phase_spatial_distribution}

La figure~\ref{fig:individual_galaxies} présente les cartes de densité de surface de l'\HI\ total, cinématiquement froid et cinématiquement chaud des galaxies individuelles. Les deux composantes s'étendent sur une grande partie des disques atomiques cartographiés et présentent des condensations et des structures allongées au sein de l'émission étendue.

La composante cinématiquement froide représente en médiane environ 45\% de la masse d'\HI\ cartographiée, contre 55\% pour la composante chaude. À l'intérieur du disque optique ($r<r_{25}$), la composante chaude domine le bilan de masse d'\HI, même si la composante froide en constitue une part importante.

La composante froide présente un aspect plus filamentaire que la composante chaude : des structures fines et allongées se détachent de l'émission atomique plus diffuse. Cela conduit à se demander si le gaz froid se trouve préférentiellement dans les bras spiraux. Une telle préférence pourrait résulter de la compression du gaz dans les bras, qui augmente sa densité et favorise son refroidissement \citep{KimSpiral2008}. Pour repérer les bras des neuf galaxies de notre échantillon présentant une structure spirale visible, nous suivons la méthode de \citet{Foyle2010}, fondée sur l'émission à 3.6 $\mu$m pour tracer les bras stellaires. Nous comparons ensuite la fraction étroite dans les bras et les régions interbras. Cette fraction n'est pas systématiquement plus élevée dans les bras.

Les modèles d'équilibre thermique de l'\HI\ prévoient la coexistence d'un gaz froid et dense et d'un gaz tiède et diffus sur un intervalle de pressions thermiques \citep{Wolfire2003}. À pression comparable, le gaz froid est plus dense en raison de sa température plus basse. Les simulations montrent qu'une augmentation de densité du gaz atomique tiède peut permettre son refroidissement et sa condensation en structures froides \citep{Saury2014}. Le gaz froid peut alors fournir de la matière aux nuages moléculaires \citep{Krumholz2009Transition}. Le H$_2$ se forme principalement à la surface des grains de poussière, et sa formation est plus rapide aux fortes densités \citep{GloverMacLow2007}. Le rayonnement ultraviolet peut toutefois dissocier les molécules. Le gaz moléculaire s'accumule lorsque la poussière et le H$_2$ lui-même absorbent suffisamment de rayonnement pour le protéger. La transition atomique--moléculaire dépend ainsi de la densité du gaz, de la densité de colonne assurant cette protection et de l'abondance de poussière \citep{Krumholz2009Transition,Sternberg2014}. Dans les simulations de \citet{GloverMacLow2007}, le gaz initialement atomique se refroidit rapidement et la compression turbulente crée des régions denses où le H$_2$ se forme plus vite. Le lien entre gaz atomique froid et formation des nuages moléculaires repose donc sur les fortes densités qu'il peut atteindre, associées à une protection suffisante. Cela motive la recherche d'une association entre la composante étroite de l'\HI\ et l'émission du CO.

Nous comparons pour cela la fraction étroite dans des régions de même aire, nettement brillantes ou faibles en CO, à des rayons galactocentriques similaires, dans huit galaxies dont l'émission CO permet ce test. Nous partons des cartes natives d'\HI\ et de CO et les convoluons pour obtenir le même faisceau. Les résolutions spatiales résultantes vont de 234 à 721 pc.

Les deux composantes atomiques peuvent être renforcées près de l'émission CO brillante, mais la fraction froide n'y est pas systématiquement plus élevée. Le cas le plus net est NGC 5194 : la fraction froide moyenne pondérée par l'aire atteint environ 45\% dans les régions brillantes en CO, contre 29\% dans les régions voisines faibles en CO. Dans NGC 6946, les deux composantes augmentent près du CO brillant, tandis que la fraction froide reste voisine de 30\%. Les autres galaxies fournissent trop peu de régions appariées pour permettre une comparaison concluante. Nous ne trouvons donc pas d'indication systématique que l'\HI\ étroit soit davantage concentré près de l'émission CO que l'\HI\ large.

L'absence d'augmentation systématique de la fraction froide n'exclut pas que le gaz atomique froid alimente les nuages moléculaires. Dans les modèles de transition atomique--moléculaire, le gaz atomique demeure dans les couches qui entourent et protègent les intérieurs moléculaires, tandis que le gaz supplémentaire peut devenir moléculaire \citep{Krumholz2009Transition,Sternberg2014}. Une fois transformé en H$_2$, l'hydrogène ne contribue plus à l'émission d'\HI. Un réservoir moléculaire croissant ne s'accompagne donc pas nécessairement d'une contribution croissante de l'\HI\ étroit. De plus, le H$_2$ peut subsister là où le CO est dissocié \citep{Wolfire2010} : les régions de comparaison faibles en CO ne sont pas nécessairement dépourvues de H$_2$. Ces effets peuvent atténuer le contraste entre régions brillantes et faibles en CO. Notre comparaison ne permet donc pas d'établir si la composante étroite trace préférentiellement les régions où le gaz atomique devient moléculaire. Il reste difficile de localiser la transition \HI--H$_2$ dans notre échantillon, puisqu'une région peut contenir du H$_2$ sans émission CO détectable \citep{Wolfire2010}.
